% exp-fips203-mlkem-kem-keygen-k3 — FIPS 203 Alg.19 ML-KEM-768 KEM KeyGen
% 编译：bash scripts/xelatex-clean.sh examples/incubating/ml-kem/ml-kem-768/exp-fips203-mlkem-kem-keygen-k3/exp-fips203-mlkem-kem-keygen-k3-实现方案-customspec.tex
\documentclass[UTF8,a4paper,11pt]{ctexart}
\usepackage{amsmath,amssymb}
\usepackage{booktabs}
\usepackage{geometry}
\usepackage{hyperref}
\usepackage{longtable}
\usepackage{array}
\usepackage{seqsplit}
\geometry{margin=2.2cm}
\newcolumntype{L}[1]{>{\raggedright\arraybackslash}p{#1}}
\newcommand{\apiname}[1]{\texttt{\seqsplit{#1}}}
\newcommand{\dirname}[1]{\texttt{\seqsplit{#1}}}
\hypersetup{colorlinks=true,linkcolor=blue,urlcolor=blue}

\title{exp-fips203-mlkem-kem-keygen-k3\\FIPS 203 Alg.19 ML-KEM-768 KEM KeyGen 实现方案}
\author{ascendc 工程（ML-KEM-768 W4b / E19）}
\date{2026-07-26}

\begin{document}
\maketitle

\begin{abstract}
本文档描述 \dirname{examples/incubating/ml-kem/ml-kem-768/exp-fips203-mlkem-kem-keygen-k3/}：
在 \textbf{Atlas A2 / Ascend910B4} 上以 AscendC 实现 FIPS 203 \textbf{Algorithm 19}，即 ML-KEM-768（$k=3$）\textbf{KEM KeyGen}。

\textbf{参数锁定}：几何、I/O 与 launch 全部锁定到
\dirname{docs/specs/fips203-mlkem768-parameter-card.md} §3.3 D19：
\texttt{ek\_kem=1184B}，\texttt{dk\_kem=2400B}，SIM/生产 \textbf{2 launch}，Alg.16 尾在第二个 launch 内融合。
\textbf{行为基线}：活跃探针
\dirname{ascendc-tests/ml-kem/ml-kem-768/pass-fix-f203-alg19-kem-keygen-device-k3/}；
PKE 主体可由本参数组 E13
\dirname{examples/incubating/ml-kem/ml-kem-768/exp-fips203-mlkem-pke-keygen-k3/}
与 Alg.16 KEM tail 组合。
\textbf{计划 AI Core 数}：Launch 1 为 AIV\_ONLY \texttt{blockDim=2}；Launch 2 为 MIX \texttt{blockDim=1}（1 AIC + 2 AIV）。
\textbf{禁止}：创建 \dirname{examples/stable/ml-kem/ml-kem-768/}；从任何 \dirname{frozen/} 目录复制、改写或移植源码；沿用 k4 的 \texttt{ek=1568}/\texttt{dk\_kem=3168} 或 pad-to-4/8 几何。
\end{abstract}

\tableofcontents
\newpage

\section{工程边界}

\subsection{允许的代码来源与依赖}

\begin{longtable}{@{}L{4.8cm}L{8.4cm}@{}}
\toprule
来源 & 用途 \\
\midrule
\dirname{ascendc-tests/ml-kem/ml-kem-768/pass-fix-f203-alg19-kem-keygen-device-k3/} & 活跃 D19 k3 行为基线；可复制后自包含适配到本 exp，最终不得保留对探针目录的编译期 include/import \\
\dirname{examples/incubating/ml-kem/ml-kem-768/exp-fips203-mlkem-pke-keygen-k3/} & 活跃 E13 incubating PKE KeyGen；提供已自包含的 k3 Alg.13 主体 \\
\dirname{examples/incubating/ml-kem/ml-kem-1024/exp-fips203-mlkem-kem-keygen-k4/} & 活跃 k4 incubating KEM KeyGen 结构参考；仅用于 retarget 到 k3 锁定字节长与几何 \\
\dirname{library/shared/} & SHAKE/Keccak、公共设备侧头文件等共享库 \\
CANN / tikicpulib / CAModel & 编译、CPU 孪生与 SIM 运行环境 \\
\bottomrule
\end{longtable}

\subsection{禁止项}

\begin{itemize}
  \item 禁止从 \dirname{ascendc-tests/frozen/} 或 \dirname{examples/frozen/} 拿源码、customspec 或路线。
  \item 禁止把 $d$、$z$、$H(\mathrm{ek})$、$\hat A$、$\hat s$、$\rho$ 等中间态作为默认 I/O 落盘；debug dump 只能由显式开关启用。
  \item 禁止把 k4 的字节长、polyvec8、$\hat A[16]$ 或零垫 4/8 当作默认。
  \item 禁止为 $H(\mathrm{ek})$ 或 $z$ 单独增加第三次生产 launch。
\end{itemize}

\section{数学语义（Alg.19 $\to$ Alg.16 $\to$ Alg.13）}

\subsection{输入输出}

\begin{longtable}{@{}L{3.5cm}L{9.8cm}@{}}
\toprule
张量 / 文件 & 契约 \\
\midrule
\texttt{input/seed\_d.bin} & 4B little-endian \texttt{uint32}，由 Host 派生 FIPS 203 的 $d$；域分离串为 \texttt{exp-mlkem-f203-2s1e-k3:SEED\_D=} \\
\texttt{input/lut\_even\_stacked.bin}、\texttt{lut\_odd\_stacked.bin} & NTT Stage2 LUT；由本目录 golden 脚本生成；不是用户可见交付输入 \\
\texttt{output/ek\_kem.bin} & 1184B，等于 PKE public key：$\text{ByteEncode}_{12}(\hat t)$ 的 1152B 后拼接 $\rho$ 的 32B \\
\texttt{output/dk\_kem.bin} & 2400B，展开布局为 $\texttt{dk\_pke}(1152)\Vert\texttt{ek\_kem}(1184)\Vert H(\texttt{ek})(32)\Vert z(32)$ \\
\bottomrule
\end{longtable}

\subsection{阶段公式}

\begin{enumerate}
  \item Alg.19：从 \texttt{seed\_d.bin} 可复现派生 $d$ 与 $z$；$d$ 使用 \texttt{exp-mlkem-f203-2s1e-k3:SEED\_D=}，$z$ 使用 \texttt{exp-mlkem-f203-kem-k3:SEED\_Z=}。
  \item Alg.16：调用 Alg.13 生成 $(\mathrm{ek}_{\mathrm{PKE}},\mathrm{dk}_{\mathrm{PKE}})$。
  \item Alg.16 尾：$\mathrm{ek}_{\mathrm{KEM}}=\mathrm{ek}_{\mathrm{PKE}}$；$\mathrm{dk}_{\mathrm{KEM}}=\mathrm{dk}_{\mathrm{PKE}}\Vert\mathrm{ek}_{\mathrm{KEM}}\Vert H(\mathrm{ek}_{\mathrm{KEM}})\Vert z$，其中 $H=\mathrm{SHA3\textrm{-}256}$。
  \item Alg.13 主体复用 E13 / D13 k3 几何：$\hat A[9,256]$、$s\Vert e[6,256]$、polyvec6 NTT、InnerProduct 输出 3 行、ByteEncode12 $3\times384$B。
\end{enumerate}

\section{AscendC 实现规划}

\subsection{Launch 与 AI Core 规模}

\begin{longtable}{@{}L{2.5cm}L{3.1cm}L{7.5cm}@{}}
\toprule
Launch & 核类型 / blockDim & 数据划分与理由 \\
\midrule
1 prep & AIV\_ONLY，\texttt{blockDim=2} & 继承 E13/D13：$\hat A$ 共 9 poly 按 5+4 分给两个 AIV；$s\Vert e$ 共 6 poly 生成到真实 polyvec6 GM。选择 2 AIV 是参数卡 §3.3 D19 对 D13 prep 的锁定继承。 \\
2 compute+tail & MIX，\texttt{blockDim=1} & 1 AIC + 2 AIV；先执行 Alg.13 行 16--21，polyvec6 NTT 连续 3+3，InnerProduct 输出 2+1；随后在同一 launch 内由约定 AIV 执行 Alg.16 KEM tail，避免第三 launch。 \\
\bottomrule
\end{longtable}

\subsection{GM 几何与偏移}

\begin{longtable}{@{}L{3.8cm}L{9.3cm}@{}}
\toprule
对象 & 锁定几何 \\
\midrule
$\hat A$ & \texttt{[9,256] int32}；9 个真实 poly，禁止 pad 到 16 \\
$s\Vert e$ & \texttt{[6,256]}；polyvec6，NTT 前后都按真实 6 poly 管理 \\
NTT S0 & \texttt{[12,256] int8}，即 \texttt{[HI6, LO6]}；每个 AIV 持有完整 poly 的 hi+lo，禁止 limbsplit \\
Stage2 mat\_c & \texttt{[48,128] int32}；逻辑 $6\times4\times2$；Cube 的 M 对齐垫到 16 只是硬件对齐，不表示假 poly \\
\texttt{ek\_kem} & 1184B；等于 \texttt{ek\_pke} \\
\texttt{dk\_kem} & 2400B；偏移：\texttt{dk\_pke[0,1152)}，\texttt{ek[1152,2336)}，\texttt{H[2336,2368)}，\texttt{z[2368,2400)} \\
\bottomrule
\end{longtable}

\subsection{同步与流水}

Prep 内部使用 UB/GM 分段搬运与必要的 \apiname{PipeBarrier}；两次 launch 之间由 Host 顺序保证 GM 可见性。
Compute 内部沿用 E13/D13 已验同步：AIV 写 Stage1/Stage3 GM 后以 CrossCore flag 与 AIC 交接；AIC 写完 \texttt{mat\_c} 后 AIV 等待再做 Stage3 merge、InnerProduct 与 ByteEncode12。
Alg.16 尾必须在 ByteEncode12 与 $ek$ fuse 完成后执行；双 AIV 写共享 GM 时须先汇合，再由固定 AIV 计算 $H(\mathrm{ek})$、派生 $z$ 并拼 \texttt{dk\_kem}。
CPU 仅作为逻辑孪生；SIM 是同步与搬运验收门槛。

\section{AscendC API 表}

本实现不引入 E19 独有的新 AscendC API；下表 API 均已有索引记录，使用约束按
\dirname{library/documents/CANN-AscendC算子开发接口参考-查阅索引.md} 执行。

\begin{longtable}{@{}L{3.2cm}L{2.7cm}L{2.4cm}L{2.3cm}L{3.0cm}@{}}
\toprule
API 名 & 分类 / 文档 & Atlas A2 & 本用例 dtype & 备注 \\
\midrule
\apiname{GetBlockIdx} / \apiname{GetSubBlockIdx} & 系统变量 / 资源 & √ & 标量索引 & 区分 AIV 分片与 MIX 子核；禁止把 CPU 打印误认为多核性能 \\
\apiname{GlobalTensor} / \apiname{LocalTensor} & 基础结构 & √ & \texttt{uint8/int8/int32} & GM/UB 张量契约必须与本 spec 几何一致 \\
\apiname{DataCopy} & Memory 数据搬运 & √ & \texttt{uint8/int8/int32} & 连续块搬运；长度与地址遵守 32B 对齐约束 \\
\apiname{Duplicate} & Memory 矢量计算 & √ & \texttt{uint8/int32} & 用于 UB 清零、pad 与 tail buffer 初始化 \\
\apiname{PipeBarrier} & Pipe 同步 & √ & — & MTE/Vector/Cube 阶段交界显式同步 \\
\apiname{CrossCoreSetFlag} / \apiname{CrossCoreWaitFlag} & MIX 同步 & √ & — & AIV/AIC 通过 GM 交接时使用；CPU 过不能删 \\
\apiname{SyncAll} & 硬同步 & √ & — & KEM tail 前双 AIV 汇合；MIX 中 AIC 已返回后按 AIV-only 汇合语义使用 \\
\apiname{Mmad} / \apiname{Fixpipe} & 矩阵计算 & √ & \texttt{int8*int8->int32} & Stage2 MMAD；逻辑 M=12，硬件 M 对齐垫到 16 \\
\apiname{ShiftRight} & 矢量算术 & √ & \texttt{int32} & limb 拆分、mask low bits 与 Barrett 相关右移 \\
\apiname{Muls} / \apiname{Add} / \apiname{Sub} & 矢量算术 & √ & \texttt{int32/int16} & limb、模加减、压缩/编码辅助 \\
\apiname{Cast} & 类型转换 & √（受限） & \texttt{int32/int16/half/int8} & 遵守索引中 A2 Cast 链约束；不得假设存在 \texttt{int32->int8} 单跳 \\
\apiname{GetValue} / \apiname{SetValue} & LocalTensor 标量访问 & √ & \texttt{uint8/int32} & 仅用于 pack 与 tail 标量缺口；不得把 Compares bit 包当逐 lane 布尔 \\
\bottomrule
\end{longtable}

\subsection{评估过但未采用}

\begin{longtable}{@{}L{4.0cm}L{9.0cm}@{}}
\toprule
API / 路线 & 不采用原因 \\
\midrule
\apiname{Compares}/\apiname{GatherMask} compact 链 & Alg.19 E19 不新增 reject compact；沿用 E13/D13 已验采样路径，不新增比较掩码风险 \\
高阶 \apiname{Matmul<>} & 本路线继承手写 MMAD + 平面 mat\_c 契约；高阶 Matmul 输出布局与 MIX 交接不作为本轮变量 \\
第三 launch KEM tail & 参数卡 §3.3 D19 锁定 2 launch；尾段必须内嵌在 compute launch 后半段 \\
\bottomrule
\end{longtable}

\section{数学步骤覆盖率}

\begin{longtable}{@{}L{4.4cm}L{4.0cm}L{2.0cm}L{2.8cm}@{}}
\toprule
数学步骤 & 实现方式 / API & 覆盖 & 缺口说明 \\
\midrule
$G(d\Vert3)$ 与 XOF & shared Keccak/SHAKE + UB 流水 & 100\% 设备 & 标量 Keccak 核属本仓 shared；非 Host oracle \\
\texttt{SampleNTT} 生成 $\hat A[9]$ & AIV\_ONLY + rej 路径 + \apiname{DataCopy} & 100\% 设备 & 5+4 分片；无零垫 \\
CBD $\eta=2$ 生成 $s\Vert e[6]$ & AIV\_ONLY + CBD LUT/向量累加 & 100\% 设备 & 输出按 polyvec6 真实行数 \\
polyvec6 NTT S1--S3 & MIX AIV/AIC + \apiname{Mmad} & 100\% 设备 & NTT 内无 \apiname{Gather} \\
InnerProduct $P=3$ & Alg.11 向量路径 + AIV 2+1 & 100\% 设备 & 非整除分片显式处理第 3 行 \\
ByteEncode12 + $ek\Vert\rho$ & ByteEncode12 向量 + 尾部标量 pack & 部分向量 & pack 尾部标量缺口沿用 E13/D13 已验实现 \\
$H(ek)$ 与 $z$ 派生 & shared Keccak SHA3-256 + tail UB 拼接 & 100\% 设备 & 仅 AIV0 负责最终拼接，前置汇合保证 GM 完整 \\
\bottomrule
\end{longtable}

\section{验证计划}

\subsection{默认验收}

\begin{verbatim}
cd examples/incubating/ml-kem/ml-kem-768/exp-fips203-mlkem-kem-keygen-k3
bash run.sh -r cpu -v Ascend910B4
SIM_DIRECT=1 bash run.sh -r sim -v Ascend910B4
\end{verbatim}

期望：\texttt{ek\_kem.bin} 为 1184B，\texttt{dk\_kem.bin} 为 2400B；两者均与 golden 对拍 PASS。
SIM 通过后记录 \texttt{Total tick} 到 \texttt{qa/active\_sim\_regress\_summary.md}，并检查用例根目录无 stray \texttt{core*.dump}、\texttt{profile\_*log*.toml}。

\subsection{可选对照}

后续可补 liboqs-768 KAT；E19 本轮门禁为 CPU + \texttt{SIM\_DIRECT=1} sim。
liboqs 仅作为黑盒 I/O oracle；禁止把 liboqs 源码实现同构移植进 AscendC。

\section{产出物}

\begin{itemize}
  \item \dirname{exp-fips203-mlkem-kem-keygen-k3-实现方案-customspec.tex}
  \item \dirname{exp-fips203-mlkem-kem-keygen-k3-实现方案-customspec.pdf}
  \item 后续实现文件须与本 customspec 同目录自包含，且自研 Python/C/AscendC 写详细中文注释。
\end{itemize}

\end{document}
